\section{Супутник Sentinel-2}

У цій роботі використано дані космічної місії %
Sentinel-2 \cite{ref1,ref2}, що дають
композити у багатьох каналах.

Композит --- це сукупність каналів, які %
використовуються для візуалізації зображення.

Перший супутник Sentinel-2A цієї місії запущено у 2015 %
році, а другий, Sentinel-2B, ---
у 2016 році.  Супутники отримують дані (таб. \ref{tab:sentinel_table}) з %
просторовою роздільністю 10, 20 та 60 м у 13
спектральних каналах \cite{ref16}.  Частота знімання на %
екваторі становить приблизно п'ять днів
\cite{ref16}.

\begin{table}[H]
    \begin{center}
        \caption{Класифікація каналів композиту супутників Sentinel-2}
        \label{tab:sentinel_table}
        \begin{tabular}{|p{1.5in} | p{1.8in} |}
            \hline
            \textbf{ Номер каналу} & \textbf{ Назва каналу} \\
            \hline
            1             & Coastal aerosol             \\
            2             & Blue                        \\
            3             & Green                       \\
            4             & Red                         \\
            5, 6, 7       & Vegetation red edge         \\
            8             & NIR                         \\
            8A            & Narrow NIR                  \\
            9             & Water vapour                \\
            10            & SWIR – Cirrus               \\
            11, 12        & SWIR                        \\
            \hline
        \end{tabular}
    \end{center}
\end{table}


Таким чином, сімейство супутників Sentinel-2 надає %
достатньо інформації для досліджень у
сфері аналізу супутникових даних і моніторингу %
земного покриву.
